% ============================================================
%  SUPPLEMENT 04 — AUTOCORRELATION DEEP DIVE
%  Pearson Correlation, Stationarity, ACF, and the Four Metrics
%  Companion to: ECDR_Data_Analysis.tex
% ============================================================
\documentclass[11pt, a4paper]{article}

\usepackage{fontspec}
\setmainfont{Noto Serif}
\setsansfont{Noto Sans}
\setmonofont{Noto Sans Mono}[Scale=0.88]

\usepackage[english]{babel}
\usepackage{amsmath, amssymb, amsthm}
\usepackage{mathtools}
\usepackage[margin=2.5cm]{geometry}
\usepackage{booktabs, tabularx, array}
\usepackage{xcolor}
\newcolumntype{Y}{>{\raggedright\arraybackslash}X}
\newcolumntype{P}[1]{>{\raggedright\arraybackslash}p{#1}}
\usepackage{hyperref}
\usepackage{enumitem}
\usepackage{microtype}
\usepackage{parskip}

\definecolor{ink}{RGB}{20,20,20}
\definecolor{teal}{RGB}{0,120,120}
\definecolor{muted}{RGB}{90,90,90}

\hypersetup{colorlinks=true, linkcolor=teal, urlcolor=teal, citecolor=teal}

\theoremstyle{definition}
\newtheorem{definition}{Definition}[section]
\newtheorem{example}{Example}[section]
\newtheorem{remark}{Remark}[section]
\theoremstyle{plain}
\newtheorem{theorem}{Theorem}[section]
\newtheorem{proposition}{Proposition}[section]

\newcommand{\R}{\mathbb{R}}
\newcommand{\E}{\mathbb{E}}
\newcommand{\Cov}{\operatorname{Cov}}
\newcommand{\Var}{\operatorname{Var}}
\newcommand{\corr}{\operatorname{corr}}
\newcommand{\code}[1]{\texttt{#1}}

\usepackage{titlesec}
\titleformat{\section}{\large\bfseries\color{teal}}{}{0em}{\thesection\quad}
\titleformat{\subsection}{\normalsize\bfseries\color{ink}}{}{0em}{\thesubsection\quad}

\begin{document}

\begin{center}
  {\LARGE\bfseries\color{teal} Autocorrelation: Concepts and the Four Metrics}\\[0.5em]
  {\large\color{muted} Pearson Correlation, Stationarity, ACF, and Their Application to HAPT IMU Data}\\[1.2em]
  {\normalsize Supplement~04 to \textit{ECDR Data Analysis} \textbar{} October~2026}\\[0.3em]
  {\small\color{muted} This document is supplementary only. The main document is not modified.}
\end{center}

\vspace{1em}\hrule\vspace{1em}

\tableofcontents
\clearpage

% ============================================================
\section{Pearson Correlation: Precise Definition}

\subsection{Population Definition}

Given two random variables $A$ and $B$ with finite variances, their Pearson
correlation coefficient is:
\begin{equation}
  \rho(A, B) \;=\; \frac{\Cov(A, B)}{\sqrt{\Var(A) \cdot \Var(B)}}
  \;=\; \frac{\E[(A - \mu_A)(B - \mu_B)]}{\sigma_A \, \sigma_B}.
\end{equation}

Properties:
\begin{itemize}[leftmargin=1.5em]
  \item $\rho \in [-1, 1]$ always (Cauchy-Schwarz inequality).
  \item $\rho = 1$: perfect positive linear relationship. $\rho = -1$: perfect negative.
  \item $\rho = 0$: no linear relationship (but may still have nonlinear dependence).
  \item \textbf{Scale-invariant}: multiplying either variable by a positive constant
        does not change $\rho$. Only the \emph{direction} of covariation matters, not magnitude.
  \item \textbf{Shift-invariant}: adding any constant to either variable does not
        change $\rho$. This is why correlation measures linear relationship after
        mean-centering.
\end{itemize}

\subsection{Sample Estimator}

For two finite vectors $a, b \in \R^m$ (same length $m$), the sample Pearson
correlation is:
\begin{equation}\label{eq:sample_pearson}
  \corr(a, b) \;=\;
  \frac{\sum_{i=1}^m (a_i - \bar{a})(b_i - \bar{b})}
       {\sqrt{\sum_{i=1}^m (a_i - \bar{a})^2 \;\cdot\; \sum_{i=1}^m (b_i - \bar{b})^2}},
\end{equation}
where $\bar{a} = m^{-1} \sum_i a_i$ and $\bar{b} = m^{-1} \sum_i b_i$.

This is the formula used in the main data analysis to compute all four
autocorrelation metrics.

% ============================================================
\section{The Autocorrelation Function (ACF)}

\subsection{Definition for a Stationary Process}

A discrete time series $\{r[n]\}_{n=0}^{T-1}$ is \textbf{(weakly) stationary} if:
\begin{enumerate}[leftmargin=1.5em]
  \item Its mean is constant: $\E[r[n]] = \mu$ for all $n$.
  \item Its autocovariance depends only on the lag $k$, not on the absolute time:
        $\Cov(r[n], r[n+k]) = \gamma(k)$ for all $n$.
\end{enumerate}

The \textbf{autocorrelation function (ACF)} at lag $k$ is:
\begin{equation}
  \rho_k \;=\; \frac{\gamma(k)}{\gamma(0)} \;=\;
  \frac{\Cov(r[n], r[n+k])}{\Var(r[n])}.
\end{equation}

Note $\rho_0 = 1$ always (a signal is perfectly correlated with itself at lag 0).

\subsection{Estimation from Finite Data}

Given a single observed series of length $T$, the sample ACF at lag $k$ is:
\begin{equation}
  \hat{\rho}_k \;=\;
  \frac{\sum_{n=0}^{T-k-1} (r[n] - \bar{r})(r[n+k] - \bar{r})}
       {\sum_{n=0}^{T-1} (r[n] - \bar{r})^2}.
\end{equation}
This is equivalent to applying the sample Pearson formula~\eqref{eq:sample_pearson}
with $a = r[0:T-k]$ and $b = r[k:T]$.

% ============================================================
\section{Physical Intuition for the HAPT Sensor Series}

\subsection{Why Lag-1 Autocorrelation Is Near 1.0}

Human motion at 50\,Hz (20\,ms between samples) is smooth. The body cannot
instantaneously change acceleration; the physics of human gait imposes continuity.
Therefore, adjacent raw samples $r[n]$ and $r[n+1]$ are nearly identical.

More precisely, the Nyquist theorem tells us that human gait energy is mostly
concentrated below 25\,Hz (and typically below 10\,Hz). At 50\,Hz sampling,
neighboring samples are oversampled relative to the signal bandwidth. The lag-1
correlation is high simply because consecutive samples see almost the same
physical state.

The empirical values ($0.83 - 0.98$ across channels, Table 5 in the main doc)
are consistent with this explanation.

\subsection{Why Lag-64 Autocorrelation Drops — Especially for Gyroscope}

Lag 64 corresponds to 64/50 = 1.28 seconds — exactly the query stride $S$.
This is a physically meaningful time scale:

\begin{itemize}[leftmargin=1.5em]
  \item \textbf{Accelerometer (acc):} Gravitational acceleration creates a strong
        DC bias in \code{acc\_x} ($+0.82\,g$). Slow postural drift changes this
        bias slowly, keeping lag-64 correlation high ($0.72 - 0.82$).
  \item \textbf{Gyroscope (gyro):} Gyroscope measures angular velocity, which
        integrates to rotation angle. In a complete gait cycle ($\approx 1$ s),
        angular velocity oscillates through a full cycle. Over 1.28 s, the
        gyroscope value can be at a completely different phase of the cycle.
        Hence lag-64 autocorrelation for \code{gyro\_x}, \code{gyro\_y},
        \code{gyro\_z} drops sharply ($0.06 - 0.15$).
\end{itemize}

\begin{remark}[Critical Implication for ECDR]
  High lag-1 autocorrelation does \textbf{not} imply that window-to-window
  differences are small. The stride is 64 samples (1.28\,s), not 1 sample.
  The relevant autocorrelation for understanding ECDR reuse is at lag 64 — and
  for gyroscope channels, the signal is nearly uncorrelated at this lag.

  This is why the main document explicitly warns: ``High short-term autocorrelation
  cannot be naively extrapolated to the 1.28\,s query stride.''
\end{remark}

% ============================================================
\section{The Four Metrics and What They Each Measure}

The main document computes four distinct autocorrelation metrics to avoid
conflating different physical phenomena. Each answers a different question.

\subsection{Metric 1: Raw Lag-1 Autocorrelation}

\textbf{Vectors:} $a = r_e[0:T_e-1, c]$ and $b = r_e[1:T_e, c]$.

\textbf{Question:} How similar is each raw sample to its immediate neighbor?

\textbf{Physical meaning:} Smoothness of the sensor signal over 20\,ms.

\textbf{Empirical values:} $0.83 - 0.98$. Very high for all channels.

\textbf{ECDR relevance:} Tells us the sensor signal is band-limited and smooth.
Does \textbf{not} tell us anything about window-to-window similarity (different time scale).

\subsection{Metric 2: Raw Lag-64 Autocorrelation}

\textbf{Vectors:} $a = r_e[0:T_e-64, c]$ and $b = r_e[64:T_e, c]$.

\textbf{Question:} How similar is each raw sample to the sample 1.28 s later?

\textbf{Physical meaning:} Temporal coherence over the query stride.

\textbf{Empirical values:} Accelerometer $0.72 - 0.82$; gyroscope $0.06 - 0.15$.

\textbf{ECDR relevance:} Directly measures how much the signal changes over one
query stride. Low gyroscope lag-64 correlation means gyroscope waveform shapes
change substantially between consecutive windows.

\subsection{Metric 3: Window Positional Correlation}

\textbf{Vectors:} For window $k$ and $k+1$ at identical tensor index $i$:
$a = X_{e,k}[c, \cdot]$ and $b = X_{e,k+1}[c, \cdot]$ (128-dimensional).

\textbf{Question:} How correlated are the \emph{shapes} of two consecutive windows
at matching temporal positions?

\textbf{Physical meaning:} Whether the waveform profile (centered and scaled) at
time offset $i$ in window $k$ matches the same offset in window $k+1$.
Note: offset $i$ in window $k$ corresponds to time $t_k + i/50$, while the same
offset in window $k+1$ is at time $t_k + (S + i)/50 = t_k + i/50 + 1.28$ s.
These are samples 1.28 s apart — Metric 3 compares the \emph{positional structure},
not the instantaneous value.

\textbf{Empirical values:} Near zero ($-0.03$ to $+0.05$). Appears surprising but is
consistent: since the stride is 1.28 s, positions $i$ in consecutive windows contain
samples 1.28 s apart, at different points in the gait cycle.

\textbf{ECDR relevance:} Near-zero positional correlation is \emph{not} a sign that
windows are unrelated. It shows that the waveform shape (as a vector aligned by
position, not by phase) is essentially random between windows — which is why
phase-aware domains (wave-domain) are more appropriate than shape-matching domains.

\subsection{Metric 4: Window Channel-Mean Correlation}

\textbf{Vectors:} Let $g_k = \frac{1}{128}\sum_i X_{e,k}[c, i]$ (channel mean
of window $k$). Compare $(g_0, \ldots, g_{N_e-2})$ and $(g_1, \ldots, g_{N_e-1})$.

\textbf{Question:} How correlated are the spatial-average values of consecutive windows?

\textbf{Physical meaning:} The mean of a window over 128 samples ($= 2.56$ s of data)
represents the ``level'' of the signal. For \code{acc\_x}, this is dominated by
gravitational orientation. This level changes slowly as posture drifts.

\textbf{Empirical values:} High ($0.56 - 0.97$). Much higher than Metric 3.

\textbf{ECDR relevance:} The spatial mean is more stable than the waveform shape.
This motivates decomposing the perturbation domain into a ``mean component''
(certifiable with a tight small domain) plus a ``waveform residual'' component
(requires wider tolerance) — exactly the idea behind the wave-domain model.

% ============================================================
\section{Worked Example: Linear Ramp Counterexample}

The main document proves (via a mathematical counterexample) that high lag-1
autocorrelation does not imply small positional window delta. Here is the
construction with full details.

\subsection{Setup}

Let $r[n] = n$ for $n = 0, 1, 2, \ldots$ (a linear ramp). The raw lag-1 correlation:
\begin{align}
  a &= (0, 1, 2, \ldots, T-2), \\
  b &= (1, 2, 3, \ldots, T-1).
\end{align}
Both have mean $\bar{a} = (T-2)/2$ and $\bar{b} = T/2$. The centered vectors are
$(a - \bar{a}) = (-T/2+1, -T/2+2, \ldots, T/2-1)$ and $(b - \bar{b}) = (-T/2, -T/2+1, \ldots, T/2-1)$.
These are identical except for the boundary element. As $T \to \infty$:
\[
  \hat{\rho}_1 \to 1.
\]
The ramp has lag-1 correlation approaching 1 because the signal is perfectly smooth.

\subsection{Window Positional Delta Is Large}

Window $k$ starts at position $kS$ (stride $S = 64$). Its tensor is:
\[
  X_{k}[i] = r[kS + i] = kS + i, \quad i = 0, \ldots, 127.
\]
Window $k+1$:
\[
  X_{k+1}[i] = (k+1)S + i.
\]
The positional delta at index $i$:
\[
  X_{k+1}[i] - X_k[i] = S = 64 \quad \text{for all } i.
\]
The $L_2$ norm of the positional delta vector:
\[
  \|X_{k+1} - X_k\|_2 = \sqrt{128 \times 64^2} = 64\sqrt{128} \approx 724.1.
\]

\textbf{Conclusion:} The ramp has lag-1 correlation $\approx 1$ but window positional
delta $\approx 724$ — two orders of magnitude away from zero. The two metrics are
completely decoupled. Lag-1 autocorrelation cannot bound the window-to-window $L_2$ distance.

% ============================================================
\section{Summary Table}

\begin{center}
\small
\begin{tabularx}{\linewidth}{@{}lYrYY@{}}
\toprule
Metric & Vectors compared & Typical value & What it measures & ECDR relevance \\
\midrule
Lag-1 raw & $r[n]$ vs $r[n+1]$ & $0.83 - 0.98$ & Signal smoothness at 20 ms & Bandwidth; irrelevant to stride \\
Lag-64 raw & $r[n]$ vs $r[n+64]$ & acc: $0.72$; gyro: $0.10$ & Coherence at 1.28 s & Directly relevant to query stride \\
Window positional & $X_k[c,i]$ vs $X_{k+1}[c,i]$ & $-0.03$ to $0.05$ & Waveform shape alignment & Near-zero: phase mismatch \\
Window channel-mean & $g_k$ vs $g_{k+1}$ & $0.56 - 0.97$ & Level / orientation stability & Supports mean-component domain \\
\bottomrule
\end{tabularx}
\end{center}

\begin{thebibliography}{9}
\bibitem{main}
  ECDR Research Workspace. \textit{ECDR Data Analysis} (English edition).
  \code{ECDR\_Data\_Analysis.tex}, October~2026.
\end{thebibliography}

\end{document}
